\chapter{Three Ways to Measure}
\label{chap:three-ways-to-measure}

% Closed question: How can an instrument compare records and repeat the result
% without pretending the comparison is the thing measured?

\begin{chapteressay}

A car passes a police officer on the side of a road. The officer raises a radar
gun. In the car, the dashboard speedometer already reads 60 mph. In the
driver's phone, a map application also estimates the car's speed from GPS
positions. Three instruments are now telling one story.

They do not tell it in the same language.

The speedometer counts rotations. A magnetic sensor sees teeth pass in the
drivetrain, counts pulses, scales by wheel circumference, and divides by elapsed
time. The radar gun sends out electromagnetic waves and compares the outgoing
frequency with the reflected frequency. The GPS receiver listens for timestamped
signals from satellites, compares arrival times, solves for position, and then
compares positions across time. The dashboard, radar display, and phone may all
say the same number, but the physical routes to that number are different.

This is the chapter's example and its warning. Agreement among instruments is
powerful only because the instruments are not the same instrument. If three
copies of the same miscalibrated speedometer agree, the agreement may only
prove that the same error has been repeated. If a pulse counter, a Doppler
instrument, and a time-of-flight system agree within their resolutions, the
comparison is stronger. The same physical claim has survived three carriers.

What is being compared? Not speed itself as a thing in the road. Each
instrument compares records. The speedometer compares pulse count to clock
count under a wheel calibration. The radar gun compares emitted and returned
frequencies under an electromagnetic calibration. The GPS receiver compares
satellite times, receiver times, and successive positions under a relativistic
clock calibration. The agreement among them is a comparison of comparisons.

A measurement therefore becomes physical law only after it can be repeated and
compared without confusing the comparison with the thing measured. A ruler does
not become a table. A clock does not become time. A speedometer does not become
speed. They produce records. Those records can agree, disagree, drift, or
calibrate one another. The instrument's interface with the world is comparison,
and the discipline of this chapter is to keep that interface honest.

The same structure appears in a kitchen scale and a measuring cup. Suppose a
baker measures flour by weight on a digital scale and by volume in a cup. The
scale compares force against an internal calibration. The cup compares a filled
volume against a marked line. The two instruments do not carry the same physical
event. Weight changes with density and packing; volume changes with how the
flour settles. The baker can still learn from their agreement or disagreement.
If the recipe says 120 grams and one cup, and the packed flour weighs 160 grams,
the disagreement is not a metaphysical crisis. It is the instrument telling the
baker that two comparison interfaces have diverged.

This chapter overlays that ordinary fact on the speedometer thread. Chapter 1
gave us admissible marks: events carried under calibration. Chapter 2 asks how
such marks can be compared. The answer must do two things at once. First, it
must let instruments repeat their readings. A comparison that cannot be repeated
is only an accident. Second, it must prevent the reading from being confused
with the world. A comparison is an interface, not an identity.

The closed question is:

\begin{center}
\emph{How can an instrument compare records and repeat the result without
pretending the comparison is the thing measured?}
\end{center}

The speedometer, radar gun, and GPS receiver answer together. Each produces a
different physical record. Each reduces its record to a ratio or comparison
under calibration. The three readings become comparable only because the
calibrations are named and the units are shared. Their agreement is not magic.
It is a constructed common interface.

By the end of the chapter, the instrument will have three new obligations. It
must expose a binary comparison: this reading is less than, equal to, or greater
than that one within the instrument's resolution. It must be repeatable: the
same physical situation, under the same calibration, must return the same
comparison. And it must remain humble: the comparison is the measurement's
admissible output, not the thing measured in itself.

\end{chapteressay}

% Phenomenon arcs use: 1/3 setup -> phenomenon box -> 2/3 structural interpretation.

\section{The Speedometer: Pulse Counting}
\label{se:the-speedometer-pulse-counting}

The speedometer's comparison begins with a pulse rate. Chapter 1 treated each
Hall-effect pulse as an admissible mark. Chapter 2 asks what can be done with a
sequence of those marks.

Suppose the sensor admits $N$ pulses in a time interval $\Delta t$. Suppose the
tone wheel has $P$ pulses per wheel revolution and the calibrated wheel
circumference is $C$. Then the instrument computes

\[
v = \frac{N}{\Delta t}\cdot\frac{C}{P}.
\]

This equation is not yet a law of nature. It is the speedometer's comparison
interface. The first ratio, $N/\Delta t$, compares pulse count to clock count.
The second ratio, $C/P$, compares distance per wheel revolution to pulses per
wheel revolution. The displayed speed is the result of composing those
comparisons.

The car's motion enters only through admitted pulses. Tire slip, vibration,
wind, road grade, and engine noise do not enter unless they change the pulse
record or its calibration. The speedometer is therefore narrow by design. It
does not measure every physical fact about the moving car. It measures the
ratio it can carry.

This narrowness is what makes comparison possible. A pulse count can be
compared with another pulse count. A clock interval can be compared with another
clock interval. A calibrated circumference can be compared with a shop standard.
The final display can be compared with a radar gun, a GPS receiver, a roadside
timing trap, or a dynamometer. The reading is useful because it is made of
comparisons all the way down.

The speedometer also shows why a reading is never merely a number. The same
number 60 can mean different things depending on the calibration. Sixty miles
per hour and sixty kilometers per hour are not the same speed. Sixty pulses per
second is not a speed until the pulses have been scaled. Sixty on a broken
display may be neither speed nor reliable frequency. The number becomes a
measurement only when the comparison grammar is present.

This gives an old physical insight a precise place in the argument.

\begin{phenom}{Archimedes-Proust Effect}
\label{ph:archimedes-proust-effect}

\PhStatement A physical measurement becomes stable when it is expressed as an
admissible ratio: one count or quantity compared with another under a named
standard.

\PhOrigin Archimedes made physical proportion operational in mechanics and
hydrostatics; Joseph Proust's law of definite proportions showed that chemical
compounds combine in fixed ratios by mass.

\PhObservation The speedometer reading is not a raw pulse count. It is a ratio
of pulses to time, scaled by the ratio of wheel circumference to pulses per
revolution. The same grammar appears in density, concentration, gear ratio, and
chemical proportion.

\PhConstraint The instrument may compare only quantities it can carry under a
calibration. A ratio with an unnamed denominator is not a physical reading.

\PhConsequence Physical comparison begins as proportion. The speedometer's
number is admissible because it is a ratio of carried records, not because the
instrument touches speed directly.

\end{phenom}

The Archimedes-Proust Effect is not a historical ornament here. It fixes the
chapter's grammar. To measure is to compare records under proportion. The
speedometer's pulse-counting interface is the simplest physical version of that
grammar in this volume.

\section{The Radar Gun: Doppler Shift}
\label{se:the-radar-gun-doppler-shift}

The radar gun agrees with the speedometer without counting wheel rotations. It
does not know the tire radius. It does not care whether the car is front-wheel
drive, rear-wheel drive, or coasting in neutral. It sends a wave and compares
what comes back.

A typical police radar gun operates in the K-band at $24{.}125\,\mathrm{GHz}$
or in the Ka-band at $33{.}4$ to $36\,\mathrm{GHz}$. The emitter, a Gunn
diode oscillator stabilized against a dielectric resonator, generates
continuous-wave microwave radiation at a precisely known frequency $f_0$. The
beam is collimated by a horn antenna to a divergence angle of about $12$
degrees. When this beam encounters a target moving along the line of sight,
the reflected wave returns with a shifted frequency. For speeds much smaller
than the speed of light, the round-trip Doppler shift is approximately
\[
\Delta f \approx \frac{2v}{c} f_0,
\]
where $v$ is the car's speed along the radar beam and $c$ is the speed of
light. The factor of two appears because the wave is shifted once as it meets
the moving car and again as it is reflected back to the receiver.

The derivation makes the chapter's point explicit. The wave's relationship to
the target involves two frequency comparisons. First, the moving car sees an
incoming wave at frequency $f_1 = f_0 (1 + v/c)$ (the target's frame sees a
blue-shifted wave because the source is approaching it). Second, the wave
reflected back to the gun is doubly shifted: the moving car acts as a
re-emitter at frequency $f_1$, but from the gun's stationary frame, the
target is moving toward the gun, so the received frequency is
$f_2 = f_1 (1 + v/c) = f_0 (1 + v/c)^2 \approx f_0 (1 + 2v/c)$ for $v \ll c$.
The shift $\Delta f = f_2 - f_0 \approx 2 f_0 v / c$ is the gun's measurable
output.

For a car at $60\,\mathrm{mph}$ ($v \approx 27\,\mathrm{m/s}$) and a K-band
radar gun at $f_0 = 24{.}125\,\mathrm{GHz}$, the predicted shift is
$\Delta f \approx 2 \times 24{.}125 \times 10^9 \times 27 / (3 \times 10^8)
\approx 4{.}3\,\mathrm{kHz}$. The radar gun extracts this audio-frequency
shift from the radio-frequency carrier by mixing the received signal with
the local oscillator output, producing a beat note at the Doppler frequency.
A simple frequency counter on the beat note gives the speed. For
police-radar accuracy of about $\pm 1\,\mathrm{mph}$, the beat frequency must
be measured to about $\pm 70\,\mathrm{Hz}$ over a one-second gating interval
--- easily achievable with standard electronics.

The instrument does not count wheel events. It compares frequencies. The
comparison is the chapter's structural point: the radar gun and the
speedometer measure the same physical quantity (the car's speed relative to
the road) through completely different instrumental procedures. The
speedometer counts mechanical rotations and multiplies by a calibrated
circumference. The radar gun emits a wave and measures a Doppler shift.
Neither instrument touches the speed directly; each performs a comparison
that, under its respective calibration, produces a number that is interpreted
as the speed.

The agreement between the two instruments, when both are properly calibrated,
is the chapter's \texttt{REPEATABLE} certificate. Disagreement is diagnostic.
If a speedometer reads 60 and the radar gun reads 55, the disagreement points
to a specific failure mode in one or both instruments. The speedometer's
characteristic failures are tire-related: a worn or oversized tire changes
the effective wheel circumference, and the calibration constant no longer
matches the physical geometry. The radar gun's characteristic failures are
geometry-related: a beam not aligned with the car's direction of motion
produces a cosine-of-angle underread; two cars in the beam force a target
selection ambiguity; oscillator drift biases the frequency comparison.

These are not reasons to distrust either instrument. They are reminders that
every comparison interface has its own failure modes, and the discipline of
measurement is to know each instrument's vulnerabilities. A radar reading is
not speed in itself; it is a Doppler comparison interpreted as speed under
named assumptions about geometry and calibration.

The radar gun also makes the communication structure of measurement explicit.
The emitted wave leaves the instrument, interacts with the car, and returns. The
reading depends on what survives that round trip. This is a channel problem:
signal, carrier, noise, encoding, decoding, and comparison.

\begin{phenom}{Fessenden-Shannon Effect}
\label{ph:fessenden-shannon-effect}

\PhStatement A measurement that travels through a channel is constrained by the
channel's ability to carry a distinguishable signal from source to receiver.

\PhOrigin Reginald Fessenden's early radio work made continuous-wave
communication operational; Claude Shannon later gave the mathematical theory of
information transmission through noisy channels.

\PhObservation A radar gun emits a calibrated microwave signal and receives a
shifted reflection. The speed reading depends on whether the returned signal can
be distinguished from noise, assigned to the correct target, and compared with
the emitted reference.

\PhConstraint The instrument may not infer more information than the channel
carries. Beam angle, noise, bandwidth, and target ambiguity limit what the
returned signal can support.

\PhConsequence Communication is already inside measurement. The radar gun is a
physical comparator because the channel preserves enough structure for emitted
and returned frequencies to be compared.

\end{phenom}

This phenomenon is the radar section's hinge. The speedometer compares events
inside the car. The radar gun compares a transmitted and returned signal across
space. In both cases, comparison requires a carried record. But the radar gun
shows that the carrier can leave the instrument, encounter the measured object,
and return with a changed structure.

That changed structure is not arbitrary. The Doppler shift is repeatable because
waves with the same emitted frequency, reflecting from targets with the same
radial speed under the same geometry, return with the same frequency shift. The
comparison is physical because the channel law is stable. The instrument's
reading is credible because the comparison can be repeated.

This is the first glimpse of a theme that will grow in later chapters. A
measurement is not only a mark in a ledger. It is also an interaction that
transports information. Chapter 7 will make that transport explicit. Here, we
need only the first lesson: the radar gun measures by comparing what it sent
with what returned, not by touching the car's speed.

\section{The GPS: Time of Flight}
\label{se:the-gps-time-of-flight}

The GPS receiver agrees with the speedometer and radar gun by using neither
wheel rotations nor reflected microwaves from the car. It measures position
from time.

The Global Positioning System consists of approximately $24$ active
satellites in six orbital planes, each at an altitude of $20{,}200\,\mathrm{km}$.
The orbital period is $11\,\mathrm{h}\,58\,\mathrm{min}$, so each satellite
completes two orbits per sidereal day and returns to roughly the same
position relative to the Earth's surface daily. The constellation geometry
is chosen so that at least four satellites are visible from any point on
the Earth's surface at any time. Each satellite carries multiple cesium and
rubidium atomic clocks with frequency stability of $10^{-13}$ over a day,
providing a precise time reference.

Each GPS satellite broadcasts a signal at two L-band frequencies:
$L1 = 1{.}57542\,\mathrm{GHz}$ and $L2 = 1{.}2276\,\mathrm{GHz}$. The
signal contains the satellite's identification code, its precise orbital
ephemeris (the Keplerian elements plus correction terms), and the time at
which the signal was sent, all encoded by phase-modulating the carrier
with pseudo-random noise codes that allow the receiver to identify each
satellite and to measure signal arrival time to within nanoseconds.

The receiver compares the transmission time encoded in each satellite's
signal with its own reception time. Multiplying the time difference by
the speed of light gives a distance to the satellite, more precisely a
\emph{pseudorange} because the receiver clock is not perfectly
synchronized with the satellite clocks. Four satellites are needed in
ordinary operation: three to triangulate the three spatial coordinates,
and a fourth to solve for the receiver clock offset (which appears as a
common bias in all four pseudoranges).

The triangulation equations are
\[
(x - x_i)^2 + (y - y_i)^2 + (z - z_i)^2 = c^2 (t_i - t + b)^2,
\quad i = 1, 2, 3, 4,
\]
where $(x, y, z)$ is the receiver's unknown position, $(x_i, y_i, z_i)$ is
the $i$-th satellite's known position at signal emission, $t_i$ is the
satellite's transmission time (broadcast in the signal), $t$ is the
receiver's reception time (read from the receiver's clock), and $b$ is
the receiver's clock offset against GPS system time. Four equations in
four unknowns ($x, y, z, b$) yield a unique solution under typical
geometric conditions.

Once the receiver has solved for its position, repeated solutions at
successive reception times yield a velocity estimate. The speed is
\[
v \approx \frac{|\Delta \mathbf{x}|}{\Delta t},
\]
where $\Delta \mathbf{x}$ is the position change between two solutions
and $\Delta t$ is the time elapsed between them.

In practice, GPS receivers use a more direct velocity measurement based
on Doppler shifts of the satellite signals themselves, which gives
better short-term velocity accuracy than position differencing. The
satellite signals arrive Doppler-shifted by both the satellite's motion
(at $3{.}87\,\mathrm{km/s}$ orbital speed) and the receiver's motion. The
receiver's contribution to the Doppler shift can be extracted from the
sum of signals from multiple satellites, yielding a velocity estimate
with accuracy on the order of $0{.}1\,\mathrm{m/s}$ even when position
accuracy is limited to a few meters.

The speedometer counts drivetrain pulses. The radar gun compares
frequencies. The GPS receiver compares times and positions across an
entire constellation of satellites. Yet all three can display the same
speed.

This agreement is possible only because the GPS receiver is aggressively
calibrated. Satellite clocks are atomic clocks. Their signals carry precise
time references. The system includes corrections for satellite orbit, clock
drift, atmospheric delay, and relativistic effects (detailed in Chapter 8).
A GPS speed reading is not ``position magic.'' It is a disciplined comparison
among many clocks.

That clock discipline brings Einstein into the chapter, but in an operational
way. Synchronization is not a philosophical luxury. If the receiver and
satellite clocks are not comparable, the position solution fails. If the clocks
are comparable but the relativistic frame corrections are wrong, the solution
drifts. The physical reading depends on an agreement about time.

\begin{phenom}{Einstein Synchronization Effect}
\label{ph:einstein-synchronization-effect}

\PhStatement Spatial measurement by light signal requires a clock-comparison
protocol. Distances inferred from signal travel time are admissible only after
the relevant clocks have been synchronized or their offsets have been solved.

\PhOrigin Einstein's 1905 analysis of simultaneity made synchronization by
light signals operational: clocks at different places are compared by a defined
signal protocol, not by access to absolute time.

\PhObservation GPS receivers solve for position and receiver-clock offset by
comparing satellite timestamps with local reception times. The speed estimate
comes only after those clock comparisons are made mutually consistent.

\PhConstraint The receiver may not treat remote timestamps as directly
comparable without a synchronization rule and calibration corrections. Time of
flight becomes distance only under that protocol.

\PhConsequence A physical instrument can compare records across space only by
carrying its reference structure with it. The GPS speed agrees with the
speedometer because its clock comparisons have been made admissible.

\end{phenom}

The GPS example closes the triangle of this chapter. The speedometer is local
and mechanical. The radar gun is remote and wave-based. The GPS receiver is
distributed and clock-based. Their agreement is not agreement of identical
apparatus. It is agreement of compatible comparisons.

That compatibility is what makes physical measurement powerful. If the
speedometer reads 60 mph, the radar gun reads 60 mph, and the GPS receiver
computes 60 mph, the car's speed is no longer tied to one fragile carrier. The
claim has survived three encodings. It has become a shared physical record.

But the humility remains. The three instruments still do not touch speed as an
independent object. They carry pulse ratios, frequency shifts, and time-of-
flight comparisons. The agreement among those records is what gives the reading
its physical force.

\section{Repeatability}
\label{se:repeatability}

Comparison without repeatability is not measurement. It is coincidence.

Suppose a car drives ten laps around a closed test track at a controlled speed.
The speedometer, radar gun, and GPS receiver are recorded on each lap. If the
track conditions, calibration, and driving profile are held fixed, the
instruments should agree within their resolutions each time. Small deviations
are not shocking; they are the measurement noise. But if the readings wander
without pattern, the comparison has failed.

Repeatability does not mean identical numbers under all conditions. It means
that the same admissible situation returns the same admissible comparison. If
the tire pressure changes, the speedometer calibration may shift. If the radar
angle changes, the Doppler reading may change. If the GPS receiver loses
satellite lock, the time-of-flight solution may degrade. Repeatability always
includes the condition ``under the same calibration and apparatus state.''

This is where a study begins. A single reading can be useful, but repeated
readings let the instrument characterize its own scatter. If the speedometer
reads 60, 60, 60, 61, 60 over five passes, the scatter is small. If the radar
gun reads 58, 60, 63, 59, 61, its scatter is larger. If the GPS receiver reads
60 steadily after filtering but lags during acceleration, that is a different
kind of repeatability limit. Repetition turns comparison into a record about the
instrument as well as about the car.

Small-sample inference belongs here because real instruments are not granted
infinite trials. A factory, laboratory, or roadside calibration procedure may
have only a handful of repeated measurements. The question is not whether
infinite repetition would settle the matter; finite instruments never get
infinite repetition. The question is what a finite repeatable record can
support.

\begin{phenom}{Gosset Effect}
\label{ph:gosset-effect}

\PhStatement Finite repeated measurements can support inference only by naming
their uncertainty. Small samples do not eliminate comparison; they force the
instrument to carry its confidence structure.

\PhOrigin William Sealy Gosset, writing as Student, developed the
small-sample $t$ distribution while working on industrial quality problems,
where only limited repeated trials were available.

\PhObservation A speedometer calibration checked over five track passes cannot
claim infinite certainty. It can report a mean discrepancy and a finite-sample
uncertainty. The comparison becomes trustworthy by carrying the uncertainty,
not by pretending the sample is unlimited.

\PhConstraint The study may not replace finite repeatability with an infinite
ideal. It must compare records using the sample actually carried by the ledger.

\PhConsequence Repeatability is physical and statistical at once. A repeated
instrument reading becomes a study only when its finite uncertainty is part of
the record.

\end{phenom}

The Gosset Effect keeps this chapter grounded. The comparison interface is not
perfect simply because it is repeated. It becomes useful when repetition exposes
the scale of its own variation. A radar gun, a speedometer, and a GPS receiver
can agree within one mile per hour; that phrase ``within'' is doing real work.
It names the tolerance that makes equality operational.

Repeatability therefore prepares the binary interface. If two readings differ
by less than the instrument's resolution, the instrument may treat them as
equal. If one reading exceeds another by more than the tolerance, it may treat
one as greater. If the difference is inside the noise band, the honest answer
may be: this instrument cannot distinguish them.

This is the beginning of order in physical comparison. Less than, equal to, and
greater than are not metaphysical verdicts. They are repeatable outcomes of a
calibrated interface.

\section{The Binary Interface}
\label{se:the-binary-interface}

The word ``binary'' here does not mean that every instrument has only two
possible readings. A speedometer has many possible displayed speeds. A GPS
receiver has many possible positions. A radar gun has many possible frequency
shifts. The binary interface is the comparison question the instrument can
answer about two records:

\[
a \leq b \quad \text{or not}.
\]

From this small question, the familiar ordering language follows. If $a \leq b$
and $b \leq a$ within the instrument's resolution, the records are equal for
that instrument. If $a \leq b$ and not $b \leq a$, then $a$ is less than $b$.
If neither comparison is admissible, the instrument may be facing incomparable
records or insufficient resolution.

For a speedometer, the binary question might be: is this reading no greater
than 60 mph? The ECU does not need to know what speed is in itself to answer.
It compares its scaled pulse ratio with the calibrated threshold. The output is
an admissible comparison. Cruise control, speed warnings, traction systems, and
data loggers all depend on such binary comparisons.

The important discipline is that the comparison is not the thing compared. A
speed warning light does not contain speed. It contains the result of a
comparison between a carried reading and a threshold. The driver's experience of
``going fast'' is also not the same as the comparison. It may correlate with
it, but the instrument's admissible output is narrower.

This is why the chapter began with three instruments. The speedometer, radar
gun, and GPS receiver can all answer comparison questions about the same car.
Is the car above 60 mph? Is its speed increasing? Is this reading less than the
previous reading? Does one instrument's output agree with another within
tolerance? The instruments answer through different carriers, but the final
interface is comparison.

There is a philosophical temptation to make the comparison disappear and say
that the instrument has simply revealed the property. This chapter resists that
temptation. The comparison is the physical interface. It is not an obstacle to
truth; it is the way a finite instrument earns truth. To pretend otherwise is to
ask the instrument to speak outside its ledger.

The binary interface also sets up the next chapter. Once an instrument can
compare recorded events, it immediately confronts what it has not recorded.
Between two pulse events, what happened? Between two GPS positions, what path
was taken? Between two radar readings, what acceleration occurred? Comparison
of records forces the question of the gap. Chapter 3 begins there.

\begin{bridge}

The chapter's question was how an instrument can compare records and repeat the
result without pretending the comparison is the thing measured.

The answer is that the instrument exposes a calibrated comparison interface.
The speedometer compares pulse count to clock count under a wheel calibration.
The radar gun compares emitted and returned frequencies through a channel. The
GPS receiver compares satellite timestamps and receiver times under a
synchronization protocol. Repetition turns those comparisons into studies, and
finite uncertainty becomes part of the record rather than an embarrassment.

The comparison remains an interface. It is powerful because it is repeatable,
portable, and checkable across instruments. It is limited because it only
compares what the instruments have carried. The next question is forced by that
limit: what can be said about what happens between recorded events without
inventing structure not in the ledger?

\end{bridge}

\begin{coda}{Three Clocks in a Railway Control Room}

In the railway control room, three clocks face the dispatcher.

The first is an old wall clock above the door. Its quartz oscillator advances a
second hand around a printed face. The second is the dispatch computer clock,
visible in the scheduling software that tracks arrivals and departures. The
third is the stationmaster's pendulum clock, maintained as a ceremonial but
working reference, its rod calibrated and its swing adjusted by hand.

The train does not touch any of them.

When the 8:10 express is due, the dispatcher compares the train's reported
position against the schedule. The wall clock says 8:08. The computer says
8:08:03. The pendulum clock says just past eight minutes after the hour. Each
clock carries time differently. The wall clock counts electronic oscillations
and gear steps. The computer counts processor-maintained time synchronized to a
network source. The pendulum counts swings under gravity, temperature, and rod
length. The clocks agree well enough for the train to be called on time.

That agreement matters because the clocks are not the same device. If the
dispatch software freezes, the wall clock continues. If the wall clock battery
weakens, the computer and pendulum expose the drift. If the pendulum is not
adjusted after a temperature change, the discrepancy appears as a comparison
among records. The control room does not possess time itself. It possesses
clocks whose readings can be compared.

The comparison is binary at the point of action. Is the train on time or not?
Is this clock ahead of that one or not? Is the departure interval less than the
minimum safe spacing or not? The dispatcher does not need a metaphysics of time
to answer the operational question. The clocks provide comparison interfaces.

Repeatability turns the control room into an instrument. If the 8:10 train
arrives at the same platform marker under the same schedule conditions, the
clocks should classify it the same way each day. If one clock disagrees
regularly, it becomes the object of calibration. If all three disagree in
different directions, the control room has lost its reference and cannot speak
honestly about the schedule.

The wall clock, computer clock, and pendulum clock perform the same structure as
the speedometer, radar gun, and GPS receiver. Each has its own carrier. Each has
its own calibration. Each can be wrong in its own way. Agreement among them is
therefore meaningful. It is not proof that any one clock touched time. It is
evidence that three different comparison chains have closed on the same
operational reading.

The train enters the record only through such comparisons. A conductor radios a
position. A track circuit reports occupancy. A platform camera confirms arrival.
The clocks do not become the train, and the schedule does not become motion.
The control room carries records and compares them. From those comparisons it
acts.

At 8:10, the express arrives. The wall clock, the computer, and the pendulum all
agree within the tolerance the railway uses. The dispatcher marks the train on
time.

That mark is not the train. It is the comparison the railway can repeat.

\end{coda}
